\documentclass{article}


% if you need to pass options to natbib, use, e.g.:
%     \PassOptionsToPackage{numbers, compress}{natbib}
% before loading neurips_2022


% ready for submission
% \usepackage{neurips_2022}
\usepackage{times}
\usepackage{latexsym}

% For proper rendering and hyphenation of words containing Latin characters (including in bib files)
\usepackage[T1]{fontenc}
\usepackage{graphicx}
\usepackage{multirow}
\usepackage[normalem]{ulem}
\useunder{\uline}{\ul}{}
\usepackage{svg}
\usepackage{hhline}

% For Vietnamese characters
% \usepackage[T5]{fontenc}
% See https://www.latex-project.org/help/documentation/encguide.pdf for other character sets

% This assumes your files are encoded as UTF8
\usepackage[utf8]{inputenc}
% \usepackage[table,xcdraw]{xcolor}
\usepackage{multirow}

% This is not strictly necessary, and may be commented out,
% but it will improve the layout of the manuscript,
% and will typically save some space.
\usepackage{microtype}
\usepackage{tabularx}
\usepackage{diagbox}

\usepackage{amsmath}

% to compile a preprint version, e.g., for submission to arXiv, add add the
% [preprint] option:
\usepackage[final]{neurips_2022}


% to compile a camera-ready version, add the [final] option, e.g.:
% \usepackage[final]{neurips_2022}


% to avoid loading the natbib package, add option nonatbib:
%    \usepackage[nonatbib]{neurips_2022}


\usepackage[utf8]{inputenc} % allow utf-8 input
\usepackage[T1]{fontenc}    % use 8-bit T1 fonts
\usepackage{hyperref}       % hyperlinks
\usepackage{url}            % simple URL typesetting
\usepackage{booktabs}       % professional-quality tables
\usepackage{amsfonts}       % blackboard math symbols
\usepackage{nicefrac}       % compact symbols for 1/2, etc.
\usepackage{microtype}      % microtypography
\usepackage{xcolor}         % colors


\title{\includegraphics[scale=0.12]{files/thinking.png} PALBERT: Teaching ALBERT to Ponder}


% The \author macro works with any number of authors. There are two commands
% used to separate the names and addresses of multiple authors: \And and \AND.
%
% Using \And between authors leaves it to LaTeX to determine where to break the
% lines. Using \AND forces a line break at that point. So, if LaTeX puts 3 of 4
% authors names on the first line, and the last on the second line, try using
% \AND instead of \And before the third author name.


\author{Nikita Balagansky, Daniil Gavrilov \\
    Tinkoff \\
  \texttt{n.n.balaganskiy@tinkoff.ai}, \texttt{d.gavrilov@tinkoff.ai}}


\begin{document}


\maketitle


\begin{abstract}
Currently, pre-trained models can be considered the default choice for a wide range of NLP tasks. Despite their SoTA results, there is practical evidence that these models may require a different number of computing layers for different input sequences, since evaluating all layers leads to overconfidence in wrong predictions (namely overthinking). This problem can potentially be solved by implementing adaptive computation time approaches, which were first designed to improve inference speed. Recently proposed PonderNet may be a promising solution for performing an early exit by treating the exit layer's index as a latent variable. However, the originally proposed exit criterion, relying on sampling from trained posterior distribution on the probability of exiting from the $i$-th layer, introduces major variance in exit layer indices, significantly reducing the resulting model's performance. In this paper, we propose improving PonderNet with a novel deterministic Q-exit criterion and a revisited model architecture. We adapted the proposed mechanism to ALBERT and RoBERTa and compared it with recent methods for performing an early exit. We observed that the proposed changes can be considered significant improvements on the original PonderNet architecture and outperform PABEE on a wide range of GLUE tasks. In addition, we also performed an in-depth ablation study of the proposed architecture to further understand Lambda layers and their performance.
\end{abstract}
\section{Introduction}

These days, fine-tuning pre-trained models on downstream tasks became a de facto standard technique for training NLP models \citep{bert, roberta, albert, gpt}. Although, as model sizes increase, it becomes harder to use them for real world applications due to computational footprint. Because of this fact, practitioners may desire to perform an early model exit (i.e., with adaptive computation time \citep{act}) to reduce the computation needed for a specific model input.



% One model that is widely used in real-world applications is ALBERT \citep{albert}, which is based on the Transformer architecture \citep{transformer} with shared layers (i.e., the same layer is evaluated several times to provide an output).

\begin{figure}[htb!]
\centering
  \includegraphics[width=\columnwidth]{files/exit2.pdf}
  \caption{A comparison of the original sampling exit criterion of PonderNet (on the left) and the proposed Q-exit criterion (on the right). PonderNet performs sampling from the Bernoulli distribution obtained from the Lambda layer at each step, possibly exiting a model too early or too late. For Q-exit, we evaluate the Cumulative distribution function (CDF) of the probability of exiting at layer $i$. Once CDF becomes greater than the threshold value ($0.5$ in this example), we perform an early exit. With such a deterministic criterion, we can perform an early exit from a model more robustly without introducing variance in the exit layer's index during inference. }
  \label{Q-exit-schema}
\end{figure}



One model that is widely used in real-world applications is ALBERT \citep{albert}, which is based on the Transformer architecture \citep{transformer} with shared layers (i.e., the same layer is evaluated several times to provide an output). Sharing layer weights allows us to think of ALBERT as of recurrent model that simplify development of early exit methods. Most of them \citep{branchynet, fastbert, deebert} rely on performing an early exit based on entropy of predictions during the model evaluation. Although, \citet{pabee} showed that running the ALBERT-Base block for a fixed number of times (10) could increase the accuracy of the fine-tuned model on specific tasks (e.g., MRPC) despite the entropy of predictions monotonously decreasing for all 12 layers. A recent PABEE solution \citep{pabee} was designed to overcome this issue by performing an early exit based on the consensus between different classifier heads from different layers. The model stops evaluating when several classifiers in a row produce the same result.

A possible orthogonal solution for this task is PonderNet \citep{pondernet} – a variational approach that treats the exit layer's index as a latent variable. By maximizing the lower bound of the likelihood of the training data, PonderNet trains a model which can predict whether it is necessary to exit from a specific layer during evaluation. However, \citet{pondernet} proposed to sample from the trained posterior distribution of exiting from each layer during inference, which leads to major variance in exit layer indices.

This paper proposes improving PonderNet adapted for ALBERT \citep{albert} and RoBERTa \citep{roberta} fine-tuning (PALBERT and PRoBERTa). Instead of performing an early exit by sampling from the trained posterior distribution during evaluation, we used a novel zero-variance exit criterion, namely \textbf{Q-exit}, which evaluates the CDF of the exit layer's probability distribution and performs a deterministic early exit. We also revisited the architectural choices of Lambda layers used to predict the probability of exiting from the current layer to make them aware of dynamics in hidden states across previous layers and the number of currently running layers.

We experimented with PALBERT and PRoBERTa on the GLUE Benchmark datasets \citep{glue}. The ablation study showed that PALBERT produced significantly better results than the original PonderNet architecture adapted for ALBERT fine-tuning. Furthermore, the proposed methods outperformed PABEE and are comparable to plain ALBERT fine-tuning, while also exceeding it in speeds. We also analyzed the trained model and provided insights on further improvement of the variational approach for early exiting\footnote{Source code: https://github.com/tinkoff-ai/palbert}.



\section{Related Work}


Most of the approaches used to perform an early exit from a model are based on the probability distribution of predictions: BranchyNet \citep{branchynet}, FastBERT \citep{fastbert}, DeeBERT \citep{deebert}, which can be seen as an entropy criterion. \citet{berxit} proposed BERxiT with LTE layer used to estimate the correctness of predictions during the inference based on the hidden state from the specific layer.

However, there is strong practical evidence that classification models' overthinking causes a reduction in predictions' entropy, making these methods difficult to use \citep{pabee}. Furthermore, it is unclear how to adapt entropy methods for regression tasks \citep{pabee}. \citet{pabee} proposed PABEE – a method to perform an early exit based on several classifiers from the different levels of a model. Once several classifiers in a row (the number of these classifiers is determined by the patience hyperparameter $t$) produce the same result, we can perform an early exit. Several recent works also utilized the idea of a consensus-based exiting strategy \citep{leebert, global}. Although, most of them mainly augment the basic consensus scheme with orthogonal heuristics. Because of this, we did not include them in our experiments and only used PABEE as a consensus-based method.

An alternative way to perform an early exit is the Ponder architecture \citep{pondernet}, which uses auxiliary Lambda layers to predict whether a model should exit from a specific layer during runtime. Inputs to Lambda layers used in PonderNet are hidden states from the current layer of a model. PonderNet can be seen as a model with the latent variable that corresponds to the exit layer index, which is trained by maximizing the lower bound of the marginalized likelihood of the data.

During inference, \citet{pondernet} proposed to sample from the trained posterior distribution of exit layer probabilities. However, this exit criterion can lead to uncertainty in exit layers for the same input (see Section \ref{prior-section}. Even if the Lambda layer produced a probability equal to $0.1$ of exiting from the first layer, we could still exit a model too early in one of ten cases, even though the probability was small. We also hypothesize that predicting exiting from a layer based entirely on a single hidden state could be sub-optimal since performing early exit could also depend on the dynamics in hidden states across layers (i.e., the Lambda layer should know how hidden states change during the evaluation).


\section{Improving PonderNet}
\label{palbert}
In this section, we describe the adaption of PonderNet to the ALBERT model (PALBERT). Although, it could be easily extended to other pre-trained models such as RoBERTa (which we also used in our experiments).


The usual ALBERT evaluation can be defined as a computation of $n$ hidden states $h_i = S(h_{i - 1})$ from the input embeddings $h_0$ of an input sequence $x$, where $i \in [1;n]$. Once $h_n$ is obtained, it is passed to a classifier block $C(h_n)$ to get the parameters of an output distribution $p(y|x)$. A common way to fine-tune this architecture on downstream tasks is to initialize the embeddings and the $S$ layer by using ALBERT (pre-trained on Masked Language Modelling) while initializing $C$ randomly and then optimizing all parameters by maximizing the likelihood of the training data.

While plain ALBERT performs a fixed number of computational steps, it is possible to perform an arbitrary number of evaluations of the layer $S$. \citet{pondernet} proposed to extend each Transformer layer with a so-called Lambda layer. More precisely, for each layer $i$, after $S$ outputs a new $h_i$, it is then passed to the classifier and Lambda layers to get parameters $C(h_i)$ of output distributions $p(y|x, i)$ and the probability of exiting from the $i$-th layer $\lambda_i = \Lambda(h_i)$, which induces a generalized geometric distribution on probability of exiting from layer $i$ equal to $p(i|x) = \lambda_i \prod_{j=1}^{i - 1}(1 - \lambda_j)$.

% \begin{equation}
% \label{posterior}
% \begin{split}
% p(i|x) = \lambda_i \prod_{j=1}^{i - 1}(1 - \lambda_j).
% \end{split}
% \end{equation}

Then, having the probability distribution from each layer $p(y|x, i)$, the parameters of the model are optimized to maximize

\begin{equation}
\label{ponder_loss}
\begin{split}
L(x, y) = \mathbb{E}_{i \sim p(i|x)}\Big[\log \big(p(y|x, i)\big)\Big] - \beta KL\Big(p(\cdot|x) || p(\cdot|\lambda)\Big) \leq \log \big(p(y|x))
\end{split}
\end{equation}

Here, $p(\cdot|\lambda)$ is a prior geometric distribution of exiting from each layer, parametrized by the hyperparameter $\lambda$, and $\mathbb{E}_{i \sim p(i|x)}\Big[\log \big(p(y|x, i)\big)\Big]$ is evaluated analytically by averaging likelihoods from different layers with posterior exit probabilities. If we treat the exit layer index as a latent variable, then optimizing $L$ from the Equation \ref{ponder_loss} could be seen as maximizing the lower bound of marginalized log likelihood $\log \big(p(y|x)\big)$ \citep{vae}.

Note that following \citet{pondernet} the prior probability of exiting from the last layer $n$ is normalized as $p(n|\lambda) = 1 - \sum_{i = 1}^{n - 1}p(i|\lambda)$ in order to make $p(i|\lambda)$ sum into $1$ with a finite number of steps, while $p(i|x)$ it truncated by the lowest layer index $j$, s.t. $\sum_{i=1}^j p(i|x) \geq 1 - 0.05$ (thus, $KL(\dots)$ is evaluated only with first $j$ layers). Also note that weights of Lambda layers are shared across layers of the model.


\subsection{Exit Criterion}


\label{q-exit}
During inference, \citet{pondernet} proposed to sample the exit layer index from $p(i|x)$ (i.e., by sampling iteratively from a Bernoulli distribution with parameter $\lambda_i$). While a sampling-based exit criterion correlates with the variational view of PonderNet's training objective, such estimation introduces the randomness in the inference process of PonderNet. 
% Since it is unpractical to sample more than one exit layer due to computational overhead, PonderNet often makes a model perform too-early or too-late exits, reducing the accuracy of the trained model.

To overcome the issue of randomness, we propose \textbf{Q-exit}\footnote{\textbf{Q}-exit stands for Quantile}: a novel deterministic criterion of performing an early exit, which we used for PALBERT. Instead of sampling from the distribution $p(i|x)$ during inference, we evaluate its CDF by accumulating $p(i|x)$ from each layer. Once the CDF is greater than the threshold hyperparameter $q$, we perform an early exit\footnote{The proposed exit criterion is the simplest deterministic criterion which we came to, since there is no ability to estimate mean or argmax statistics during the inference without running all layers of the model, which is impractical if we want to perform an early exit.}. See Figure \ref{Q-exit-schema} for a schematic comparison of the sampling criterion with Q-exit. Threshold $q$ can be seen as a trade-off between underthinking and overthinking. Therefore, $q$ should be selected during the validation of the trained model to choose the best-performing value.

Based on our experiments, we found that the proposed criterion produced significantly better accuracy on various tasks compared to the original sampling criterion (see Sections \ref{ablation-section}, \ref{glue-section}), while also being more practical than the original sampling criterion.


\subsection{Lambda Layer Architecture}
While the original PonderNet used a single layer MLP to obtain logit of exiting probability, we hypothesize that making the Lambda layer understand the dynamics of changing ALBERT hidden states is crucial for achieving good performance. To do so, instead of passing a single hidden state $h_i$ from the $i$-th layer in $\Lambda$, we concatenate it with $h_{i - 1}$. I.e., for PALBERT, we evaluate the probability of exiting from $i$-th layer as $\lambda_i = \Lambda([h_i, h_{i - 1}])$.

% \begin{equation}
% \label{new_lambda}
% \lambda_i = \Lambda([h_i, h_{i - 1}]).
% \end{equation}

We used a $3$ layer MLP with $\tanh$ activation for the Lambda layer to operate with more complex input. Based on the ablation study, we observed that increasing the capacity improves the accuracy of the trained model (See section \ref{ablation-section}). We also found it beneficial to fine-tune the Lambda layer with a different learning rate than all other parameters.


\begin{table*}
\centering
\resizebox{0.8\columnwidth}{!}{%
\begin{tabular}{ccccccc}
\toprule
\multicolumn{4}{c|}{Method} & \multicolumn{1}{c|}{SST-2} & \multicolumn{1}{c|}{RTE} & CoLA \\ \midrule
\multicolumn{4}{c|}{ALBERT} & \multicolumn{1}{c|}{92.7   ±   0.3} & \multicolumn{1}{c|}{77.0 ± 1.9} & 57.0   ±   2.1 \\ \midrule
\multicolumn{4}{c|}{ALBERT + PonderNet} & \multicolumn{1}{c|}{91.1   ±   0.6} & \multicolumn{1}{c|}{73.5 ± 1.9} & 50.8   ±   2.2 \\ 

\multicolumn{4}{c|}{ALBERT + PonderNet (Closed-Form Expectation)} & \multicolumn{1}{c|}{92.3 ± 0.4} & \multicolumn{1}{c|}{76.8 ± 3.0} & 55.9 ± 2.2 \\ \midrule

\multicolumn{1}{c}{Q-exit} & \multicolumn{1}{c}{Lambda LR} & \multicolumn{1}{c}{3-Layer Lambda} & $h$ concat. & \multicolumn{3}{c}{} \\ \midrule
\multicolumn{1}{c}{+} & \multicolumn{1}{c}{-} & \multicolumn{1}{c}{-} & - & \multicolumn{1}{|c|}{92.2   ±   0.3} & \multicolumn{1}{c|}{77.3 ± 1.4} & 55.7   ±   0.9 \\ \midrule
\multicolumn{1}{c}{+} & \multicolumn{1}{c}{+} & \multicolumn{1}{c}{-} & - & \multicolumn{1}{|c|}{92.7   ±   0.4} & \multicolumn{1}{c|}{77.3 ± 1.4} & 56.5   ±   1.2 \\ \midrule
\multicolumn{1}{c}{+} & \multicolumn{1}{c}{+} & \multicolumn{1}{c}{+} & - & \multicolumn{1}{|c|}{92.6   ±   0.3} & \multicolumn{1}{c|}{77.0 ± 1.4} & 56.3   ±   2.4 \\ \midrule
\multicolumn{1}{c}{+} & \multicolumn{1}{c}{+} & \multicolumn{1}{c}{-} & + & \multicolumn{1}{|c|}{\textbf{93.0   ±   0.3}} & \multicolumn{1}{c|}{76.5 ± 1.6} & 56.9   ±   1.9 \\ \midrule
\multicolumn{1}{c}{+} & \multicolumn{1}{c}{+} & \multicolumn{1}{c}{+} & + & \multicolumn{1}{|c|}{92.9   ±   0.2} & \multicolumn{1}{c|}{\textbf{77.8 ± 1.2}} & \textbf{57.4   ±   1.7} \\ \bottomrule

\end{tabular}
}
\caption{An ablation study of the proposed PALBERT architecture. "Lambda LR" corresponds to fine-tuning the Lambda layer with its learning rate, "3-layer Lambda" refers to making the Lambda layer have three MLP layers instead of one, and "$h$ concat." stands for concatenation of two hidden states as input to the Lambda layer.}
  \label{ablation-study-table}
\end{table*}


\section{Experiments}
\label{experiments}
\subsection{Ablation Study}
\label{ablation-section}
We performed an ablation study of the proposed changes in PonderNet architecture with PALBERT. 

We experimented with adding the proposed Q-exit criterion, Lambda layer architecture, and fine-tuning strategies. These methods were benchmarked on SST-2, RTE, and CoLA tasks from the GLUE Benchmarking dataset \citep{glue}.

We compared proposed changes with plain ALBERT and PonderNet adapted for ALBERT. Also, in these experiments we compared with PonderNet evaluated in expectation of predictions (e.g., evaluated $p(y|x) = \mathbb{E}_{i \sim p(i|x)}\Big[p(y|x, i)\Big]$ in closed-form). While such a model is impractical for early exiting setup (it does not perform any early exit but evaluates all layers to estimate the expectation), it could show a gap between closed-form expectation of model predictions and its single sample Monte Carlo estimation.

For evaluation, we performed a grid hyperparameter search on an appropriate metric score on the dev split for each dataset. Following the PABEE training setup, we trained all models with a fixed learning rate until validation metrics stopped increasing for $5$ epochs. We used Adam optimizer \citep{adam} for all experiments, a fixed $q=0.5$ on models with the Q-exit criterion, as well as a fixed classifier dropout value equal to $0.1$ \citep{dropout}, and $\lambda = 0.1$.

We trained each model $5$ times with the best hyperparameters and reported the mean and std values. A full list of the methods' hyperparameter ranges can be found in Table \ref{ablation-hyp-range}.

\begin{table*}
\centering
\resizebox{\columnwidth}{!}{%
\begin{tabular}{c|cccccccc|c} 
\toprule
Method & SST-2 & RTE & QNLI & CoLA & MRPC & MNLI & QQP & STS-B & Macro \\ 
\midrule
\multicolumn{10}{c}{Dev set} \\ 
\midrule
ALBERT & 92.7 & 76.5 & 91.5 & 56.6 & 90.5 & 84.8 & 88.9 & 90.6 & 84.0 \\
ALBERT (9 L.) & 91.1 & 74.2 & 91.2 & 56.9 & 89.7 & 84.3 & 89.7 & 90.1 & 83.4 \\ 
\midrule
ALBERT PABEE & 92.7 & 76.9 & \textbf{91.5} & 55.6 & 88.3 & 84.5 & \textbf{88.9} & \textbf{89.9} & 83.5 \\
ALBERT PonderNet & 91.3 & 74.0 & 88.3 & 51.3 & 87.1 & 81.7 & 87.7 & 88.2 & 81.2 \\
PALBERT (ours) & \textbf{93.1} & \textbf{78.3} & 91.0 & \textbf{58.1} & \textbf{89.3} & \textbf{ 84.7} & \textbf{88.9} & \textbf{89.9} & \textbf{84.2} \\ 
\midrule
\multicolumn{10}{c}{Test set} \\ 
\midrule
ALBERT & 93.4 & 70.0 & 92.1 & 50.5 & 85.6 & 79.0 & 84.7 & 87.4 & 80.3 \\ 
\midrule
ALBERT PABEE & 92.7 & 71.1 & 91.3 & 46.0 & 84.3 & 79.2 & 83.7 & \textbf{86.5} & 79.3 \\
ALBERT PonderNet & 90.2 & 68.6 & 88.5 & 43.7 & 84.3 & 81.1 & 77.4 & 83.6 & 77.2 \\
PALBERT (ours) & \textbf{93.0} & \textbf{73.5} & \textbf{91.7} & \textbf{48.6} & \textbf{87.1} & \textbf{79.8} & \textbf{84.4} & \textbf{86.5} & \textbf{80.6} \\
\bottomrule
\end{tabular}}
\caption{A comparison of PALBERT with recent approaches on the GLUE benchmark.  In the Macro column, we present the average results across tasks. We bolded the best results. Note that we did not bold the ALBERT rows since there is no early exit applied and we reported it for reference.}
  \label{glue-dev-albert}
\end{table*}


\begin{table*}
\vskip 0.1in
\centering
\resizebox{\columnwidth}{!}{%
\begin{tabular}{c|cccccccc|c} 
\toprule
Method & SST-2 & RTE & QNLI & CoLA & MRPC & MNLI & QQP & STS-B & Macro \\ 
\midrule
\multicolumn{10}{c}{Dev set} \\ 
\midrule
RoBERTa & 94.5~ & 79.4 & 92.4 & 63.1 & 91.3 & 86.7~ & 89.8 & ~90.7 & 86.0 \\ 
\midrule
RoBERTa PABEE & \textbf{93.8} & 76.9 & \textbf{91.9} & 60.7 & 89.6 & 86.4 & \textbf{91.1} & \textbf{89.9} & 85.0 \\
PRoBERTa (ours) & 93.7 & \textbf{79.1} & 91.6 & \textbf{62.2} & \textbf{90.2} & \textbf{86.7} & \textbf{91.1} & \textbf{89.9} & \textbf{85.6} \\ 
\midrule
\multicolumn{10}{c}{Test set} \\ 
\midrule
RoBERTa & 94.5 & 74.1 & 93.1 & 58.9 & 89.1 & 86.5 & 80.9 & 88.4 & 83.2 \\ 
\midrule
RoBERTa PABEE & \textbf{95.4} & 71.7 & \textbf{92.0} & 51.2 & 87.6 & \textbf{86.2} & \textbf{80.4} & \textbf{86.2} & 81.3 \\
PRoBERTa (ours) & 94.5 & \textbf{72.1} & 91.8 & \textbf{58.6} & \textbf{88.3} & 86.1 & 79.6 & 85.7 & \textbf{82.1} \\
\bottomrule
\end{tabular}}
\caption{A comparison of PRoBERTa with recent approaches on the GLUE benchmark.  In the Macro column, we present the average results across tasks. We bolded the best results. Note that we did not bold the RoBERTa rows since there is no early exit applied and we reported it for reference.}
\label{glue-dev-roberta}
\end{table*}

See Table \ref{ablation-study-table} for the full list of the results of our ablation study. Based on these experiments, PonderNet architecture is seen as performing worse than vanilla ALBERT fine-tuning. Also, PonderNet with closed-form evaluation of predictions expectation performed better than the sampling criterion, indicating that randomness in exit layers leads to poor performance. 

At the same time, the deterministic Q-exit criterion improves PonderNet accuracy when compared to a random sampling of the exit layer and is comparable to the evaluation of PonderNet in expectation. A more complex Lambda layer that can handle the dynamics of hidden state changes can further improve model accuracy when compared to the original PonderNet.

% \newpage
\subsection{GLUE Experiments}
\label{glue-section}

We compared PALBERT with PonderNet architecture adapted for ALBERT fine-tuning, and PABEE with a fixed patience value $t = 6$ on all GLUE tasks. We also trained Ponder RoBERTa (PRoBERTa) and compared it with PABEE and adapted for RoBERTa. Note that we re-implemented all baselines used in our experiment for consistency in comparison. See Appendix Section \ref{reimplement-section} for the details of reimplementation.


Following the experimental setup from Section \ref{ablation-section}, we trained $5$ models with the best hyperparameters across the hyperparameter search and reported the median task score on the dev set. We evaluated the test scores on the best models, selected based on their dev scores. We report the two metrics' mean for the MRPC, QQP, and STS-B tasks. For the MNLI task, we report the mean accuracy across matched and mismatched datasets. 

See Tables \ref{glue-dev-albert},\ref{glue-dev-roberta} for the full list of results. We observed that PALBERT outperformed PABEE on a wide range of tasks. Vanilla PonderNet with the sampling exit criterion performed the worst. Vanilla ALBERT outperformed PABEE on most tasks and is comparable to PALBERT, while the latter has the highest score averaged across all tasks (see Macro column in Tables \ref{glue-dev-albert}, \ref{glue-dev-roberta}). 

For experiments with RoBERTa, we observed that PRoBERTa outperformed PABEE on Dev split, while it either performed better or marginally worse for the Test set. Thus, PRoBERTa showed a higher macro score compared to PABEE.

Note that PABEE is performing poorly for tasks with a small dataset (e.g., CoLA, RTE). We hypothesize that this is caused by several independent classifiers at each layer $C_i$ failing to train well enough, whereas PALBERT was capable of utilizing knowledge sharing between layers.



\begin{table}[htb!]
\centering
\resizebox{0.4\columnwidth}{!}{%
\begin{tabular}{l|l} 
\toprule
Parameter                  & Values range                \\ 
\midrule
Learning rate              & {[}1e-5, 2e-5, 3e-5, 5e-5]  \\ 
\midrule
Batch size                 & {[}16, 32, 128]             \\ 
\midrule
Lambda learning rate       & {[}1e-5, 2e-5, 3e-5]        \\ 
\midrule
$\beta$   & {[}0.5]                     \\ 
\bottomrule
\end{tabular}}
\caption{Hyperparameter search ranges used in all of our experiments. }
  \label{ablation-hyp-range}
\end{table}


\begin{figure}[htb!]
 \centering
 
  \includegraphics[width=0.49\columnwidth]{files/sst2.pdf}
  \includegraphics[width=0.49\columnwidth]{files/rte.pdf}
  \includegraphics[width=0.49\columnwidth]{files/cola.pdf}
  \includegraphics[width=0.49\columnwidth]{files/mrpc.pdf}
  
  \caption{PALBERT score dependency on the Q-exit threshold. We report the mean and std values of task metrics across $5$ trained models. See section \ref{q-understanding} for more details}
  \label{threshold-acc-fig}
\end{figure}



\subsection{Understanding the Threshold of Q-exit}
\label{q-understanding}
As noted previously in Section \ref{q-exit}, we treat the threshold value $q$ of the Q-exit criterion as a trade-off between underthinking and overthinking, where increasing $q$ forces a model to evaluate more layers, and vice versa.

Therefore, it is necessary to find the best-performing threshold for each task where a model has the highest accuracy. To do so, we evaluated trained PALBERT models from the ablation study (see Section \ref{ablation-section}) on dev splits of tasks with different values of $q$. We then averaged obtained metrics and reported the mean and std values for various thresholds (See Figure \ref{threshold-acc-fig} for the results). 

We observed that exiting models with $q=0.5$ shows the best overall performance for different tasks. Making $q$ greater than $0.5$ leads to a reduction in accuracy and can often force models to evaluate all $12$ layers of ALBERT-Base.

% We associate such behavior of trained models with the fact that the huge probability mass of trained posterior probability $p(i|x)$ is concentrated near the last layers of models (see Figure \ref{sampling-probs-fig}).

% We hypothesize that the reason for this is that the parameterization of prior probability $p(i|\lambda)$ as geometric distribution with normalized last layer, proposed with PonderNet \citep{pondernet}, leads to a huge prior probability of exiting on the last layers (See Section \ref{palbert}). For MRPC, we observe a huge variance in the probabilities of exiting from different models on the first layers, which we believe leads to poor performance on this task. 

% Note that these plots could be seen as an estimation of probabilities of exiting from each layer with vanilla PonderNet sampling exit criterion. For the RTE task, layers $i \in [1; 10]$ have approximately the same probability of exiting with a total probability mass close to $0.5$, introducing huge variance in model outputs. 
It is also notable that PALBERT, having a large threshold value $q$ that performs constant exit on the last layer, has better accuracy than vanilla ALBERT fine-tuned for the SST-2 task. This indicates that adding auxiliary tasks for each layer of ALBERT improves performance compared to plain fine-tuning with a single classifier on the last layer of the model.


% \section{On the Correlation with Complexity of Samples}
% We hypothesize that the index of the exit layer should be correlated with the complexity of the input sequence.

\subsection{Speed Analysis}
\label{speedup-section}
While making $q<0.5$ improves inference speed, it can also lead to underthinking and lower accuracy (see Figure \ref{speedup-accuracy-fig}). We compared PALBERT using different threshold values $q$ to PABEE with different patience values $t$, which stands for the number of layers necessary to output the same result in a row to perform an early exit. We evaluated task scores for the specified hyperparameters as well as the increase in speed when compared to vanilla ALBERT inference of a full model with $12$ layers.

Overall, we observed that PALBERT mostly produced higher scores on different tasks while also being slightly faster than PABEE. For the CoLA and MRPC datasets, PALBERT performed better. The proposed method outperformed PABEE while achieving the same increase in speed.

We observed questionable results on the SST-2 dataset: the best score for the PABEE model is slightly higher than that of PALBERT. However, it was obtained with a negligible increase in speed compared to ALBERT, as the best-performing patience for this setup is $11$ layers while the whole model has only $12$ layers.

Furthermore, unlike PALBERT, PABEE performed significantly worse than plain ALBERT fine-tuned on the CoLA and RTE tasks.


\begin{figure}
 \centering
 \includegraphics[width=0.49\columnwidth]{files/speedup_sst2.pdf}
 \includegraphics[width=0.49\columnwidth]{files/speedup_rte.pdf}
 \includegraphics[width=0.49\columnwidth]{files/speedup_cola.pdf}
 \includegraphics[width=0.49\columnwidth]{files/speedup_mrpc.pdf}
  \caption{A comparison between PALBERT and PABEE models on SST-2, RTE, CoLA, and MRPC tasks. We varied the threshold value of Q-exit for PALBERT and the patience hyperparameter for PABEE to obtain the plots of task scores of inference increasing in speed. 1x stands for plain ALBERT inference without performing an early exit. The horizontal line corresponds to plain ALBERT fine-tuning. See Section \ref{speedup-section} for the analysis of these plots.}
  \label{speedup-accuracy-fig}
\end{figure}



% \section{Results}
% The proposed deterministic Q-exit criterion significantly outperformed the vanilla PonderNet sampling criterion. Changes in the architecture of Lambda layers further improved obtained models (see Section \ref{ablation-section}). While the proposed Q-exit criterion has the parameter $q$, which should be selected for each task, we observed that $q = 0.5$ generally performed best across different tasks (see Section \ref{q-understanding}). 

% Based on all performed experiments, we conducted that PALBERT in general performed better than related baselines such as PABEE (see Tables \ref{glue-dev-albert}, \ref{glue-dev-roberta}), and is comparable with vanilla ALBERT fine-tuning while being faster than it. We also observed that PALBERT outperforms PABEE on small datasets, while PABEE can reach slightly better results with tasks with a huge amount of training data (see Section \ref{glue-section}). 

% We also noted that performing constant exit on the last layer produced better results than vanilla ALBERT fine-tuning for some tasks (e.g., SST-2). Based on these results, we hypothesize that reducing over-thinking is only a part of increasing the accuracy of trained models. One could also concentrate on improving the fine-tuning process of a pre-trained model (e.g., add auxiliary tasks on each Transformer layer).




\begin{figure}[htb!]
 \centering

  \includegraphics[width=0.49\columnwidth] {files/mrpc_exit_hist.pdf}
  \includegraphics[width=0.49\columnwidth] {files/rte_exit_hist.pdf}
    \includegraphics[width=0.49\columnwidth] {files/mrpc_seeds.pdf}
      \includegraphics[width=0.49\columnwidth] {files/rte_seeds.pdf}

  \caption{Histogram of exit layer indices with Q-exit criterion and different prior distribution parameters $\lambda$ (top) alongside with task metrics for trained models (bottom) for MRPC and RTE tasks. See Section \ref{prior-section} for more details.}
  \label{lambda-hist-fig}
\end{figure}



\begin{figure}[htb!]
 \centering
 
  \includegraphics[width=0.23\columnwidth] {files/sst_density.pdf}
  \includegraphics[width=0.23\columnwidth] {files/rte_density.pdf}
  \includegraphics[width=0.23\columnwidth] {files/cola_density.pdf}
  \includegraphics[width=0.23\columnwidth] {files/mrpc_density.pdf}
  
  \caption{An estimation of $\mathbb{E}_{x \sim D}\Big[p(i|x)\Big]$, where $p(i|x)$ is a trained posterior probability of exiting from layer $i$ of PALBERT models across different tasks, $\lambda = 0.1$, and $D$ is the distribution of the training dataset. We took $5$ models trained on these tasks and sampled exit layer indices for the training dataset's inputs. We smoothed the obtained probabilities for visibility. Note that these probabilities could be seen as the distribution of exit layers across different tasks for plain PonderNet with sampling criterion, which indicates on high variance in indices of output layers.}
  \label{sampling-probs-fig}
\end{figure}

% \newpage
\subsection{On the Prior of PALBERT}
\label{prior-section}

While we used a fixed value of $\lambda = 0.1$ in our main experiment following \citet{pondernet} (see Section \ref{glue-section}), we also investigate the effect of changing the parameter of the prior distribution on exit layer indices.

We followed the experimental setup from the previous sections and trained models with $\lambda \in {[}0.1, 0.15, 0.25, 0.5]$. For each $\lambda$, we had its own best hyperparameter set, which was then used to train $5$ models.

See Figure \ref{lambda-hist-fig} for the results. We observed that varying $\lambda$ affects the performance of the trained model since it has a huge impact on the distribution of exit layer indices. Although the best metric for the RTE task was reached with $\lambda = 0.08$, such a model mostly performs an exit from the last layer. Meanwhile, $\lambda = 0.08$ performs marginally worse than $\lambda = 0.1$ for MRPC.

Note that while Figure \ref{lambda-hist-fig} was obtained using the Q-exit criterion, we could also estimate $\mathbb{E}_{x \sim D}\Big[p(i|x)\Big]$ (see Figure \ref{sampling-probs-fig}), which is the distribution of exit layer indices with vanilla sampling criterion proposed by \citet{pondernet}. Even if $\lambda$ is low enough so that the mode of distribution is the $12$-th layer, in almost half of the cases, the sampling criterion forces the model to exit from one of the first $10$ layers. We hypothesize that such randomness in the index of the exit layer leads to the poor performance observed during previous experiments (see Section \ref{glue-section}).

From such perspective, dependency on the explicit prior distribution could be seen as the main limitation of the proposed method, since treating the exit layer index as a latent variable justified its prior distribution. If proper prior distribution was found, then the well-performing model could be trained, and vice versa. 


\section{Future Work and Limitations}
\label{conclusion}

In this paper, we proposed improving the PonderNet architecture to perform an early exit using fine-tuned ALBERT and RoBERTa models with the novel Q-exit criterion and a revisited Lambda layer architecture. This approach is orthogonal to the consensus-based approaches \citep{pabee, leebert} highly represented in the field.

While PALBERT and PRoBERTa outperformed some recent State-of-The-Art methods used for an early exit, there is a clear direction for further improvement of this method, as it was not capable of outperforming plain ALBERT and RoBERTa on some GLUE tasks.

We believe that PALBERT could benefit from the development of a new parameterization of the prior distribution on exiting from each layer since it directly affects the resulting posterior distribution used to perform an early exit (see Section \ref{prior-section}). However, the development of an appropriate prior distribution for early exiting models is orthogonal to the development of a proper exit criterion.

Furthermore, there is no theoretical justification for the Q-exit threshold value. Although we observed that $q=0.5$ performed best, it is without a clear explanation as to why that is so. We hypothesize that bringing more insights into developing deterministic exit criteria could further improve the proposed method.

In addition, adding more auxiliary tasks could also make it possible to improve PALBERT further. This way, PALBERT training can be made more PABEE-like by making independent classifiers for each layer of the model or adding self-distillation across layers.

While the variational view of the early exiting mechanism allows us to train models which could be exited early, we observed that the overall performance of the trained model highly depends on the prior distribution of exit layer indices (see Section \ref{prior-section}). Because of this, the proposed method requires more computation for hyperparameter search compared to PABEE \citep{pabee}. We believe that, in order to further develop training-level early exiting, it may be beneficial to move away from variational inference to some other form of trainable mechanism without relying on exit layers' explicit prior knowledge.



\bibliography{custom}
\bibliographystyle{acl_natbib}


%%%%%%%%%%%%%%%%%%%%%%%%%%%%%%%%%%%%%%%%%%%%%%%%%%%%%%%%%%%%


\appendix
\section{Re-implementation of Baselines}
\label{reimplement-section}
While \citet{pabee} provided official implementation of the proposed method, we have several concerns regarding its reproducibility and experimental setup. The original implementation implies running evaluation iteration-wise (i.e., the best performing model is selected once in $n$ iterations). 

In this case, several iterations $n$ between validation epochs could be seen as an additional hyperparameter, which has a large impact on the resulting performance. Although, \citet{pabee} did not include the best performing value so that the results could be easily reproduced. 

Because of this, we decided to re-implement PABEE strictly following the original paper and training details, while evaluating the model after every epoch, instead of evaluating after $n$ training steps, which we believe is a fairer way to compare methods.


% \section*{Checklist}


% \begin{enumerate}


% \item For all authors...
% \begin{enumerate}
%   \item Do the main claims made in the abstract and introduction accurately reflect the paper's contributions and scope?
%     \answerYes{See Section 4.}
%   \item Did you describe the limitations of your work?
%     \answerYes{See Section 5}
%   \item Did you discuss any potential negative societal impacts of your work?
%     \answerNA{We proposed a method that allows us to perform an early exit from an ALBERT model. We believe that there are no any potential negative societal impacts of our work.}
%   \item Have you read the ethics review guidelines and ensured that your paper conforms to them?
%     \answerYes{}
% \end{enumerate}


% \item If you are including theoretical results...
% \begin{enumerate}
%   \item Did you state the full set of assumptions of all theoretical results?
%     \answerNA{}
%         \item Did you include complete proofs of all theoretical results?
%     \answerNA{}
% \end{enumerate}


% \item If you ran experiments...
% \begin{enumerate}
%   \item Did you include the code, data, and instructions needed to reproduce the main experimental results (either in the supplemental material or as a URL)?
%     \answerYes{We ran all experiments on openly available GLUE benchmark datasets and provide source code for our paper. We also provide WandB configs for hyperparameter search and best model evaluation on dev split of the data so that the results could be easily reproduced.}
%   \item Did you specify all the training details (e.g., data splits, hyperparameters, how they were chosen)?
%     \answerYes{We conducted all experiments on default splits of the GLUE benchmarking datasets. Hyperparameter ranges are available in Table 4.}
%         \item Did you report error bars (e.g., with respect to the random seed after running experiments multiple times)?
%     \answerYes{We reported mean and std values across different seeds for the ablation study (see Table 1), while for the main experiment, we decided to report only median values across runs for consistency with recent works. During the exploration of the proposed method in Sections 4.3, 4.5, we also reported std values across runs as well as mean values.}
%         \item Did you include the total amount of computing and the type of resources used (e.g., type of GPUs, internal cluster, or cloud provider)?
%     \answerNo{We did not include this information due to NDA concerns.}
% \end{enumerate}


% \item If you are using existing assets (e.g., code, data, models) or curating/releasing new assets...
% \begin{enumerate}
%   \item If your work uses existing assets, did you cite the creators?
%     \answerYes{See Section 4.1}
%   \item Did you mention the license of the assets?
%     \answerNA{}{}
%   \item Did you include any new assets either in the supplemental material or as a URL?
%     \answerNA{}
%   \item Did you discuss whether and how consent was obtained from people whose data you're using/curating?
%     \answerNA{}
%   \item Did you discuss whether the data you are using/curating contains personally identifiable information or offensive content?
%     \answerNA{}
% \end{enumerate}


% \item If you used crowdsourcing or conducted research with human subjects...
% \begin{enumerate}
%   \item Did you include the full text of instructions given to participants and screenshots, if applicable?
%     \answerNA{}
%   \item Did you describe any potential participant risks, with links to Institutional Review Board (IRB) approvals, if applicable?
%     \answerNA{}
%   \item Did you include the estimated hourly wage paid to participants and the total amount spent on participant compensation?
%     \answerNA{}
% \end{enumerate}


% \end{enumerate}
 
% 
\end{document}